\chapter{Laplace's Equation}

\section*{2.1. Mean Values}

\begin{theorem}[Mean value theorem]\label{gt.thm2.1}\dcref{ch2:2.1}\group{chapter-2}
\source{gilbarg-trudinger:page-0026}
Let $u\in C^2(\Omega)$ and let $B_R(y)\Subset\Omega$. If
$\Delta u=0$ (respectively $\Delta u\geq0$ or $\Delta u\leq0$), then
\begin{equation}
\tag{2.5}
u(y)=(\leq,\ \geq)\frac{1}{n\omega_nR^{n-1}}
\int_{\partial B_R(y)}u\,ds
\end{equation}
and
\begin{equation}
\tag{2.6}
u(y)=(\leq,\ \geq)\frac{1}{\omega_nR^n}
\int_{B_R(y)}u\,dx.
\end{equation}
\end{theorem}
\begin{proof}
  \uses{gt.thm2.1}
\sketch The divergence theorem shows that the derivative with respect to
$\rho$ of the spherical mean over $\partial B_\rho(y)$ has the sign of
$\Delta u$. Compare this mean with its limit $u(y)$ as $\rho\downarrow0$ to
obtain (2.5), then integrate the spherical inequality in $\rho$ to obtain
(2.6).
\end{proof}

\section*{2.2. Maximum and Minimum Principles}

\begin{theorem}[Strong maximum principle]\label{gt.thm2.2}\dcref{ch2:2.2}\group{chapter-2}
\source{gilbarg-trudinger:page-0027}
\uses{gt.thm2.1}
If $\Delta u\geq0$ in a domain $\Omega$ and $u$ attains its supremum at an
interior point, then $u$ is constant. The analogous assertion holds for
$\Delta u\leq0$ and an interior minimum. Thus a nonconstant harmonic
function has no interior extremum.
\end{theorem}
\begin{proof}\uses{gt.thm2.1}
\sketch The set where $u$ equals its maximum is relatively closed. The ball
mean inequality from Theorem~\ref{gt.thm2.1} shows that this set contains a
neighborhood of each of its points, hence is also relatively open. Connectedness
finishes the proof; replace $u$ by $-u$ for minima.
\end{proof}

\begin{theorem}[Weak maximum principle]\label{gt.thm2.3}\dcref{ch2:2.3}\group{chapter-2}
\source{gilbarg-trudinger:page-0027}
Let $\Omega$ be bounded and
$u\in C^2(\Omega)\cap C^0(\overline\Omega)$. Then
\begin{equation}
\tag{2.7}
\Delta u\geq0\Longrightarrow
\sup_\Omega u=\sup_{\partial\Omega}u,
\qquad
\Delta u\leq0\Longrightarrow
\inf_\Omega u=\inf_{\partial\Omega}u.
\end{equation}
In particular, a harmonic function lies between its boundary extrema.
\end{theorem}
\begin{proof}
  \uses{gt.thm2.3}
\sketch An extremum exists on the compact closure. If it were achieved only
in the interior, Theorem~\ref{gt.thm2.2} would make $u$ constant.
\end{proof}

\begin{theorem}[Dirichlet uniqueness]\label{gt.thm2.4}\dcref{ch2:2.4}\group{chapter-2}
\source{gilbarg-trudinger:page-0027}
\uses{gt.thm2.3}
If $u,v\in C^2(\Omega)\cap C^0(\overline\Omega)$ satisfy
$\Delta u=\Delta v$ in bounded $\Omega$ and $u=v$ on $\partial\Omega$, then
$u=v$ in $\Omega$.
\end{theorem}
\begin{proof}
  \uses{gt.thm2.3, gt.thm2.1}
\sketch Apply Theorem~\ref{gt.thm2.3} to the harmonic function $u-v$ and
to its negative.
\end{proof}

\section*{2.3. Harnack's Inequality}

\begin{theorem}[Harnack inequality]\label{gt.thm2.5}\dcref{ch2:2.5}\group{chapter-2}
\source{gilbarg-trudinger:page-0028}
\uses{gt.thm2.1}
If $u\geq0$ is harmonic in $\Omega$ and
$\Omega'\Subset\Omega$ is bounded, then
\begin{equation}
\tag{2.8}
\sup_{\Omega'}u\leq C\inf_{\Omega'}u,
\end{equation}
where $C$ depends only on $n$, $\Omega'$, and $\Omega$.
\end{theorem}
\begin{proof}
  \uses{gt.thm2.1, gt.thm2.5}
\sketch For $x_1,x_2\in B_R(y)$ with $B_{4R}(y)\subset\Omega$, the two ball
mean formulas give
\begin{equation}
\tag{2.9}
u(x_1)\leq3^nu(x_2).
\end{equation}
Cover an arc in $\Omega$ joining
points of maximum and minimum on $\overline\Omega'$ by finitely many such
balls and iterate this inequality.
\end{proof}

\section*{2.4. Green's Representation}

For $u,v\in C^2(\overline\Omega)$ on a domain where the divergence theorem
holds, Green's identities are
\begin{equation}
\tag{2.10}
\int_\Omega(v\Delta u+Du\mathbin{\cdot}Dv)\,dx
 =\int_{\partial\Omega}v\,\partial_\nu u\,ds.
\end{equation}
\begin{equation}
\tag{2.11}
\int_\Omega(v\Delta u-u\Delta v)\,dx
 =\int_{\partial\Omega}(v\partial_\nu u-u\partial_\nu v)\,ds.
\end{equation}
The normalized fundamental solution is
\begin{equation}
\tag{2.12}
\Gamma(x-y)=
\begin{cases}
\dfrac{|x-y|^{2-n}}{n(2-n)\omega_n},&n>2,\\
\dfrac{1}{2\pi}\log|x-y|,&n=2.
\end{cases}
\end{equation}
For $n>2$,
\[
D_i\Gamma(x-y)=\frac{(x_i-y_i)|x-y|^{-n}}{n\omega_n},
\]
and
\begin{equation}
\tag{2.13}
D_{ij}\Gamma(x-y)=\frac{|x-y|^2\delta_{ij}-n(x_i-y_i)(x_j-y_j)}{n\omega_n}|x-y|^{-n-2}.
\end{equation}
Thus $\Gamma$ is harmonic away from $y$ and
\begin{equation}
\tag{2.14}
|D_i\Gamma(x-y)|\leq\frac{|x-y|^{1-n}}{n\omega_n},
\qquad
|D_{ij}\Gamma(x-y)|\leq\frac{|x-y|^{-n}}{\omega_n},
\end{equation}
more generally $|D^\beta\Gamma(x-y)|\leq
C(n,|\beta|)|x-y|^{2-n-|\beta|}$. On
$\Omega\setminus\overline{B_\rho(y)}$, (2.11) reads
\begin{equation}
\tag{2.15}
\int_{\Omega-B_\rho}\Gamma\Delta u\,dx
=\int_{\partial\Omega}(\Gamma\partial_\nu u-u\partial_\nu\Gamma)\,ds
+\int_{\partial B_\rho}(\Gamma\partial_\nu u-u\partial_\nu\Gamma)\,ds.
\end{equation}
Sending $\rho\downarrow0$ gives
\begin{equation}
\tag{2.16}
u(y)=\int_{\partial\Omega}
\left(u\partial_\nu\Gamma-\Gamma\partial_\nu u\right)\,ds
+\int_\Omega\Gamma\Delta u\,dx.
\end{equation}
For compactly supported $u$ this reduces to
\begin{equation}
\tag{2.17}
u(y)=\int_\Omega\Gamma(x-y)\Delta u(x)\,dx,
\end{equation}
while harmonic $u$ satisfies
\begin{equation}
\tag{2.18}
u(y)=\int_{\partial\Omega}
(u\partial_\nu\Gamma-\Gamma\partial_\nu u)\,ds.
\end{equation}

If $h\in C^1(\overline\Omega)\cap C^2(\Omega)$ is harmonic, then
\begin{equation}
\tag{2.19}
-\int_{\partial\Omega}(u\partial_\nu h-h\partial_\nu u)\,ds
=\int_\Omega h\Delta u\,dx.
\end{equation}
Consequently $G(x,y)=\Gamma(x-y)+h(x,y)$ satisfies
\begin{equation}
\tag{2.20}
u(y)=\int_{\partial\Omega}(u\partial_\nu G-G\partial_\nu u)\,ds
+\int_\Omega G\Delta u\,dx.
\end{equation}
If $G(\cdot,y)=0$ on $\partial\Omega$, it is the Dirichlet Green function and
\begin{equation}
\tag{2.21}
u(y)=\int_{\partial\Omega}u\partial_\nu G\,ds
+\int_\Omega G\Delta u\,dx.
\end{equation}
The representation shows that harmonic functions are analytic.

\section*{2.5. The Poisson Integral}

For $B_R=B_R(0)$, define the inverse point by
\begin{equation}
\tag{2.22}
\bar x=\frac{R^2}{|x|^2}x\quad(x\ne0),
\end{equation}
with $\bar0=\infty$. The Green function is
\begin{equation}
\tag{2.23}
G(x,y)=
\begin{cases}
\Gamma(|x-y|)-\Gamma\left(\dfrac{|y|}{R}|x-\bar y|\right),&y\neq0,\\
\Gamma(|x|)-\Gamma(R),&y=0.
\end{cases}
\end{equation}
It satisfies
\begin{equation}
\tag{2.24}
G(x,y)=G(y,x),\qquad G(x,y)\leq0
\quad(x,y\in\overline B_R),
\end{equation}
and its outer normal derivative at $x\in\partial B_R$ is
\begin{equation}
\tag{2.25}
\partial_\nu G(x,y)=
\frac{R^2-|y|^2}{n\omega_nR}|x-y|^{-n}.
\end{equation}
Consequently a harmonic $u\in C^0(\overline B_R)$ satisfies the Poisson
formula
\begin{equation}
\tag{2.26}
u(y)=\frac{R^2-|y|^2}{n\omega_nR}
\int_{\partial B_R}\frac{u(x)}{|x-y|^n}\,ds_x.
\end{equation}

\begin{theorem}[Dirichlet problem on a ball]\label{gt.thm2.6}\dcref{ch2:2.6}\group{chapter-2}
\source{gilbarg-trudinger:page-0032}
\source{gilbarg-trudinger:page-0033}
\uses{gt.thm2.5}
For continuous $\varphi$ on $\partial B_R$, define
\begin{equation}
\tag{2.27}
u(x)=\frac{R^2-|x|^2}{n\omega_nR}
\int_{\partial B_R}\frac{\varphi(y)}{|x-y|^n}\,ds_y
\quad(x\in B_R),
\end{equation}
and set $u=\varphi$ on $\partial B_R$. Then
$u\in C^2(B_R)\cap C^0(\overline B_R)$ and $\Delta u=0$ in $B_R$.
\end{theorem}
\begin{proof}
  \uses{gt.thm2.5, gt.thm2.6, gt.thm2.2, gt.thm2.3, gt.thm2.4}
\sketch The Poisson kernel satisfies
\begin{equation}
\tag{2.28}
\int_{\partial B_R}K(x,y)\,ds_y=1\qquad(x\in B_R),
\end{equation}
where
\begin{equation}
\tag{2.29}
K(x,y)=\frac{R^2-|x|^2}{n\omega_nR|x-y|^n}
\qquad(x\in B_R,\ y\in\partial B_R).
\end{equation}
It is harmonic and nonnegative. Near $x_0\in\partial B_R$, split the integral into a
small cap, controlled by continuity of $\varphi$, and its complement, where
the kernel tends uniformly to zero. Thus $u(x)\to\varphi(x_0)$.
\end{proof}

\section*{2.6. Convergence Theorems}

\begin{theorem}[Mean value characterization]\label{gt.thm2.7}\dcref{ch2:2.7}\group{chapter-2}
\source{gilbarg-trudinger:page-0033}
\uses{gt.thm2.6, gt.thm2.2, gt.thm2.3, gt.thm2.4}
A continuous function $u$ is harmonic in $\Omega$ if and only if every ball
$B_R(y)\Subset\Omega$ satisfies
\begin{equation}
\tag{2.30}
u(y)=\frac{1}{n\omega_nR^{n-1}}\int_{\partial B_R(y)}u\,ds.
\end{equation}
\end{theorem}
\begin{proof}\uses{gt.thm2.6, gt.thm2.2, gt.thm2.3, gt.thm2.4}
\sketch Necessity is Theorem~\ref{gt.thm2.1}. Conversely, solve the
Dirichlet problem on each ball using Theorem~\ref{gt.thm2.6}. The difference
between $u$ and that harmonic solution has the mean value property and zero
boundary values, so the maximum-principle argument forces it to vanish.
\end{proof}

\begin{theorem}\label{gt.thm2.8}\dcref{ch2:2.8}\group{chapter-2}
\source{gilbarg-trudinger:page-0033}
A uniform limit of harmonic functions is harmonic.
\end{theorem}
\begin{proof}
  \uses{gt.thm2.5, gt.thm2.8}
\sketch Pass to the limit in the spherical mean identity and apply
Theorem~\ref{gt.thm2.7}.
\end{proof}

\begin{theorem}[Harnack convergence theorem]\label{gt.thm2.9}\dcref{ch2:2.9}\group{chapter-2}
\source{gilbarg-trudinger:page-0034}
\uses{gt.thm2.5, gt.thm2.8}
Let $u_j$ increase monotonically through harmonic functions on $\Omega$. If
$u_j(y)$ is bounded at one point, then $u_j$ converges uniformly on every
bounded $\Omega'\Subset\Omega$ to a harmonic function.
\end{theorem}
\begin{proof}
  \uses{gt.thm2.5, gt.thm2.8, gt.thm2.9}
\sketch Apply Harnack's inequality to the nonnegative harmonic differences
$u_m-u_j$. Their values at $y$ are Cauchy, hence the differences tend
uniformly to zero on compact subdomains. Theorem~\ref{gt.thm2.8} identifies
the limit as harmonic.
\end{proof}

\section*{2.7. Interior Derivative Estimates}

The mean value and divergence theorems yield, for harmonic $u$,
\begin{equation}
\tag{2.31}
|Du(y)|\leq\frac{n}{d_y}\sup_\Omega|u|,
\qquad d_y=\dist(y,\partial\Omega).
\end{equation}

\begin{theorem}[Higher interior derivative estimates]
\label{gt.thm2.10}
\dcref{ch2:2.10}
\group{chapter-2}
\source{gilbarg-trudinger:page-0035}
If $u$ is harmonic in $\Omega$ and $\Omega'\Subset\Omega$, then for every
multi-index $\alpha$,
\begin{equation}
\tag{2.32}
\sup_{\Omega'}|D^\alpha u|
\leq\left(\frac{n|\alpha|}{d}\right)^{|\alpha|}
\sup_\Omega|u|,
\qquad d=\dist(\Omega',\partial\Omega).
\end{equation}
\end{theorem}
\begin{proof}
\sketch Apply (2.31) successively to derivatives of $u$ on
$|\alpha|$ equally spaced nested subdomains.
\end{proof}

\begin{theorem}[Compactness of bounded harmonic families]
\label{gt.thm2.11}
\dcref{ch2:2.11}
\group{chapter-2}
\source{gilbarg-trudinger:page-0035}
Every bounded sequence of harmonic functions on $\Omega$ has a subsequence
converging uniformly on compact subdomains to a harmonic function.
\end{theorem}
\begin{proof}
\sketch Estimate (2.32) gives local equiboundedness and equicontinuity.
Arzela--Ascoli and a diagonal exhaustion produce a locally uniform limit,
which is harmonic by Theorem~\ref{gt.thm2.8}.
\end{proof}

\section*{2.8. Perron's Method}

A continuous $u$ is \emph{subharmonic} in $\Omega$ if, on every
$B\Subset\Omega$, it lies below each harmonic $h$ with $u\leq h$ on
$\partial B$; superharmonicity reverses the inequalities. Subharmonic
functions satisfy the strong maximum principle, are stable under finite
maxima, and remain subharmonic under the harmonic lifting
\begin{equation}
\tag{2.33}
U=\begin{cases}
\bar u,&\text{in }B,\\
u,&\text{in }\Omega\setminus B,
\end{cases}
\end{equation}
where $\bar u$ is the harmonic extension of $u|_{\partial B}$.

For bounded boundary data $\varphi$, let $S_\varphi$ be the continuous
subharmonic functions on $\overline\Omega$ satisfying
$v\leq\varphi$ on $\partial\Omega$.

\begin{theorem}[Perron solution]\label{gt.thm2.12}\dcref{ch2:2.12}\group{chapter-2}
\source{gilbarg-trudinger:page-0036}
\source{gilbarg-trudinger:page-0037}
\uses{gt.thm2.9}
The function
$u(x)=\sup_{v\in S_\varphi}v(x)$ is harmonic in $\Omega$.
\end{theorem}
\begin{proof}
  \uses{gt.thm2.9, gt.thm2.12}
\sketch At a fixed $y$, choose subfunctions approaching $u(y)$ and replace
them by their harmonic liftings in a ball about $y$. The compactness theorem
gives a harmonic limit $v\leq u$ with $v(y)=u(y)$. If $v<u$ elsewhere,
take the maximum with another subfunction, lift again, and use the strong
maximum principle to contradict equality at $y$. Hence $u=v$ locally.
\end{proof}

A \emph{barrier} at $\xi\in\partial\Omega$ is a superharmonic
$w\in C^0(\overline\Omega)$ with $w(\xi)=0$ and
$w>0$ on $\overline\Omega\setminus\{\xi\}$. The point $\xi$ is
\emph{regular} if it has a barrier. A local barrier extends to a global one
by truncating it at a positive constant outside a small neighborhood.

\begin{lemma}\label{gt.lem2.13}\dcref{ch2:2.13}\group{chapter-2}
\source{gilbarg-trudinger:page-0037}
\source{gilbarg-trudinger:page-0038}
\uses{gt.thm2.12}
Let $u$ be the Perron solution. If $\xi$ is regular and $\varphi$ is
continuous at $\xi$, then $u(x)\to\varphi(\xi)$ as $x\to\xi$.
\end{lemma}
\begin{proof}
  \uses{gt.thm2.12, gt.lem2.13}
\sketch For a barrier $w$ and any $\varepsilon>0$, choose $k$ so that
$\varphi(\xi)\pm\varepsilon\pm kw$ are super- and subfunctions relative to
$\varphi$. Then
\[
|u(x)-\varphi(\xi)|\leq\varepsilon+kw(x),
\]
and $w(x)\to0$ at $\xi$.
\end{proof}

\begin{theorem}[Dirichlet solvability and regularity]\label{gt.thm2.14}\dcref{ch2:2.14}\group{chapter-2}
\source{gilbarg-trudinger:page-0038}
\uses{gt.lem2.13}
The classical Dirichlet problem on a bounded domain is solvable for every
continuous boundary value if and only if every boundary point is regular.
\end{theorem}
\begin{proof}\uses{gt.lem2.13}
\sketch If every point is regular, Lemma~\ref{gt.lem2.13} gives the correct
boundary values for the Perron solution. Conversely, solve the problem with
boundary data $\varphi(x)=|x-\xi|$; the solution is a barrier at $\xi$.
\end{proof}

An exterior tangent ball $\overline{B_R(y)}\cap\overline\Omega=\{\xi\}$
produces the barrier
\begin{equation}
\tag{2.34}
w(x)=\begin{cases}
R^{2-n}-|x-y|^{2-n},&n\geq3,\\
\log(|x-y|/R),&n=2.
\end{cases}
\end{equation}
Thus the exterior sphere condition, in particular a $C^2$ boundary, implies
regularity.

\section*{2.9. Capacity}

For a bounded smooth set $E\subset\mathbb R^n$, $n\geq3$, let $u$ be its
exterior conductor potential. For any smooth closed $\Sigma$ enclosing $E$,
\begin{equation}
\tag{2.35}
\operatorname{cap}E
=-\int_\Sigma\partial_\nu u\,ds
=\int_{\mathbb R^n\setminus E}|Du|^2\,dx.
\end{equation}
Equivalently,
\begin{equation}
\tag{2.36}
\operatorname{cap}E
=\inf_{\substack{v\in C_0^1(\mathbb R^n)\\v=1\text{ on }E}}
\int_{\mathbb R^n}|Dv|^2\,dx.
\end{equation}
Extend capacity to compact sets by nested smooth approximation. For
$\lambda\in(0,1)$ and $x_0\in\partial\Omega$, put
\[
C_j=\operatorname{cap}\{x\notin\Omega:|x-x_0|\leq\lambda^j\}.
\]
The Wiener criterion says that $x_0$ is regular exactly when
\begin{equation}
\tag{2.37}
\sum_{j=0}^\infty\frac{C_j}{\lambda^{j(n-2)}}=\infty.
\end{equation}
